\documentclass[10pt,a4paper]{amsart}
\usepackage{amsmath,amssymb,amsthm,graphicx}
% Retained minimal preamble for the Stacks Project chapter more-morphisms.tex (More on Morphisms).
% Verbatim excerpt of the chapter preamble: Xy-pic setup (native preamble lines 14-18) and the native
% theorem declarations with the mathematical macros (native preamble lines 154-217). Copied mechanically;
% nothing added, removed or redefined. The wrapper class and the amsmath/amssymb/amsthm/graphicx base
% packages are supplied by the assembler, so the class/externaldocument part of the native preamble is omitted.
% For commutative diagrams we use Xy-pic
\usepackage[all]{xy}

% We use 2cell for 2-commutative diagrams.
\xyoption{2cell}
% Theorem environments.
%
\theoremstyle{plain}
\newtheorem{theorem}[subsection]{Theorem}
\newtheorem{proposition}[subsection]{Proposition}
\newtheorem{lemma}[subsection]{Lemma}

\theoremstyle{definition}
\newtheorem{definition}[subsection]{Definition}
\newtheorem{example}[subsection]{Example}
\newtheorem{exercise}[subsection]{Exercise}
\newtheorem{situation}[subsection]{Situation}

\theoremstyle{remark}
\newtheorem{remark}[subsection]{Remark}
\newtheorem{remarks}[subsection]{Remarks}

\numberwithin{equation}{subsection}

% Macros
%
\def\lim{\mathop{\mathrm{lim}}\nolimits}
\def\colim{\mathop{\mathrm{colim}}\nolimits}
\def\Spec{\mathop{\mathrm{Spec}}}
\def\Hom{\mathop{\mathrm{Hom}}\nolimits}
\def\Ext{\mathop{\mathrm{Ext}}\nolimits}
\def\SheafHom{\mathop{\mathcal{H}\!\mathit{om}}\nolimits}
\def\SheafExt{\mathop{\mathcal{E}\!\mathit{xt}}\nolimits}
\def\Sch{\mathit{Sch}}
\def\Mor{\mathop{\mathrm{Mor}}\nolimits}
\def\Ob{\mathop{\mathrm{Ob}}\nolimits}
\def\Sh{\mathop{\mathit{Sh}}\nolimits}
\def\NL{\mathop{N\!L}\nolimits}
\def\CH{\mathop{\mathrm{CH}}\nolimits}
\def\proetale{{pro\text{-}\acute{e}tale}}
\def\etale{{\acute{e}tale}}
\def\QCoh{\mathit{QCoh}}
\def\Ker{\mathop{\mathrm{Ker}}}
\def\Im{\mathop{\mathrm{Im}}}
\def\Coker{\mathop{\mathrm{Coker}}}
\def\Coim{\mathop{\mathrm{Coim}}}

% Boxtimes
%
\DeclareMathSymbol{\boxtimes}{\mathbin}{AMSa}{"02}

%
% Macros for moduli stacks/spaces
%
\def\QCohstack{\mathcal{QC}\!\mathit{oh}}
\def\Cohstack{\mathcal{C}\!\mathit{oh}}
\def\Spacesstack{\mathcal{S}\!\mathit{paces}}
\def\Quotfunctor{\mathrm{Quot}}
\def\Hilbfunctor{\mathrm{Hilb}}
\def\Curvesstack{\mathcal{C}\!\mathit{urves}}
\def\Polarizedstack{\mathcal{P}\!\mathit{olarized}}
\def\Complexesstack{\mathcal{C}\!\mathit{omplexes}}
% \Pic is the operator that assigns to X its picard group, usage \Pic(X)
% \Picardstack_{X/B} denotes the Picard stack of X over B
% \Picardfunctor_{X/B} denotes the Picard functor of X over B
\def\Pic{\mathop{\mathrm{Pic}}\nolimits}
\def\Picardstack{\mathcal{P}\!\mathit{ic}}
\def\Picardfunctor{\mathrm{Pic}}
\def\Deformationcategory{\mathcal{D}\!\mathit{ef}}
\usepackage{cleveref}
\crefname{originalreference}{original reference}{original references}
\Crefname{originalreference}{Original reference}{Original references}
\makeatletter
\newlabel{algebra-lemma-constructible-is-image}{{00F8}{source}{}{sourcepointer-1}{}}
\newlabel{algebra-lemma-constructible-is-image@cref}{{[originalreference][0][]00F8}{[1][0][]source}{}{}{}}
\makeatother
\crefname{originallemma}{original lemma}{original lemmas}
\Crefname{originallemma}{Original lemma}{Original lemmas}
\makeatletter
\newlabel{lemma-quasi-finite-separated-pass-through-finite}{{05K0}{source}{}{sourcepointer-2}{}}
\newlabel{lemma-quasi-finite-separated-pass-through-finite@cref}{{[originallemma][0][]05K0}{[1][0][]source}{}{}{}}
\makeatother
\crefname{originallemma}{original lemma}{original lemmas}
\Crefname{originallemma}{Original lemma}{Original lemmas}
\makeatletter
\newlabel{lemma-weighting-base-change}{{0F3B}{source}{}{sourcepointer-3}{}}
\newlabel{lemma-weighting-base-change@cref}{{[originallemma][0][]0F3B}{[1][0][]source}{}{}{}}
\makeatother
\crefname{originallemma}{original lemma}{original lemmas}
\Crefname{originallemma}{Original lemma}{Original lemmas}
\makeatletter
\newlabel{lemma-weighting-open}{{0GK8}{source}{}{sourcepointer-4}{}}
\newlabel{lemma-weighting-open@cref}{{[originallemma][0][]0GK8}{[1][0][]source}{}{}{}}
\makeatother
\crefname{originalreference}{original reference}{original references}
\Crefname{originalreference}{Original reference}{Original references}
\makeatletter
\newlabel{morphisms-lemma-massage-finite}{{03HW}{source}{}{sourcepointer-5}{}}
\newlabel{morphisms-lemma-massage-finite@cref}{{[originalreference][0][]03HW}{[1][0][]source}{}{}{}}
\makeatother
\crefname{originalreference}{original reference}{original references}
\Crefname{originalreference}{Original reference}{Original references}
\makeatletter
\newlabel{morphisms-theorem-chevalley}{{054K}{source}{}{sourcepointer-6}{}}
\newlabel{morphisms-theorem-chevalley@cref}{{[originalreference][0][]054K}{[1][0][]source}{}{}{}}
\makeatother
\crefname{originalreference}{original reference}{original references}
\Crefname{originalreference}{Original reference}{Original references}
\makeatletter
\newlabel{properties-section-constructible}{{054B}{source}{}{sourcepointer-7}{}}
\newlabel{properties-section-constructible@cref}{{[originalreference][0][]054B}{[1][0][]source}{}{}{}}
\makeatother
\crefname{originalreference}{original reference}{original references}
\Crefname{originalreference}{Original reference}{Original references}
\makeatletter
\newlabel{topology-section-constructible}{{04ZC}{source}{}{sourcepointer-8}{}}
\newlabel{topology-section-constructible@cref}{{[originalreference][0][]04ZC}{[1][0][]source}{}{}{}}
\makeatother
\title{Level sets of a weighting are locally constructible}
\author{The Stacks Project Authors}
\date{stacks-project:more-morphisms.tex:local-snapshot-20261011; GNU FDL 1.2 or later}
\begin{document}
\maketitle
\noindent\textit{Editorial note.} Counted span is the author's own proof environment, source lines 22520-22575, immediately following the statement at 22513-22518.\par\medskip
\noindent\textit{Editorial note.} The proof contains a footnote about immersions and nonzero loci; it is retained verbatim inside the span.\par\medskip
\noindent\textit{Editorial note.} Five cross-chapter pointers (Algebra, Morphisms, Properties, Topology) are recorded in assembly.external\_references.\par\medskip
\noindent\textit{Editorial note.} Steps the author leaves to the reader are flagged here for the reviewer: source line(s) 22528 carry an "omitted"/"follows easily"/"immediate from the definitions" phrase. Those sentences are retained verbatim; the displayed argument around them stays inside the counted span.\par\medskip
\noindent\textit{Allowed standard prerequisites (editorial).}\begin{itemize}
\item Openness of the locus where a weighting is nonzero. chapter lemma lemma-weighting-open, used from proof line 22521 without further mention. Gives openness of the nonzero locus, which is the basic input for constructibility of level sets.
\item Base change of weightings. chapter lemma lemma-weighting-base-change, used from proof line 22521. Lets the argument be carried out after changing the base.
\item Local constructibility and constructible subsets. Topology, section topology-section-constructible and Properties, section properties-section-constructible, cited at proof lines 22524-22526. Supplies the closure properties of locally constructible sets used to combine the local pieces.
\item Chevalley's theorem and massaging finite morphisms. Morphisms, theorem morphisms-theorem-chevalley and lemma morphisms-lemma-massage-finite, cited in the proof. Makes images of constructible sets constructible and normalises a finite morphism into a disjoint union of finite locally free pieces.
\item Passing through finite morphisms. chapter lemma lemma-quasi-finite-separated-pass-through-finite, cited in the proof. Reduces the locally quasi-finite situation to a finite one.
\end{itemize}\medskip
\ifdefined\hypertarget\hypertarget{sourcepointer-1}{}\fi
\noindent\textit{Source reference.} Original reference 00F8: focused-seeds/geometry/S9-morph-jumps-weighting-level-sets/ref-tags.txt, source lines 4--4. Navigation pointer only: the counted proof cites this statement by \textbackslash{}ref as a standing input; the deduction the counted proof has to make is carried out inside the counted span, and no part of that deduction is delegated to this pointer.\par\medskip
\ifdefined\hypertarget\hypertarget{sourcepointer-2}{}\fi
\noindent\textit{Source reference.} Original lemma 05K0: sources/stacks-more-morphisms/pinned-original-source.tex, source lines 12457--12457. Navigation pointer only: the counted proof cites this statement by \textbackslash{}ref as a standing input; the deduction the counted proof has to make is carried out inside the counted span, and no part of that deduction is delegated to this pointer.\par\medskip
\ifdefined\hypertarget\hypertarget{sourcepointer-3}{}\fi
\noindent\textit{Source reference.} Original lemma 0F3B: sources/stacks-more-morphisms/pinned-original-source.tex, source lines 22173--22173. Navigation pointer only: the counted proof cites this statement by \textbackslash{}ref as a standing input; the deduction the counted proof has to make is carried out inside the counted span, and no part of that deduction is delegated to this pointer.\par\medskip
\ifdefined\hypertarget\hypertarget{sourcepointer-4}{}\fi
\noindent\textit{Source reference.} Original lemma 0GK8: sources/stacks-more-morphisms/pinned-original-source.tex, source lines 22239--22239. Navigation pointer only: the counted proof cites this statement by \textbackslash{}ref as a standing input; the deduction the counted proof has to make is carried out inside the counted span, and no part of that deduction is delegated to this pointer.\par\medskip
\ifdefined\hypertarget\hypertarget{sourcepointer-5}{}\fi
\noindent\textit{Source reference.} Original reference 03HW: focused-seeds/geometry/S9-morph-jumps-weighting-level-sets/ref-tags.txt, source lines 5--5. Navigation pointer only: the counted proof cites this statement by \textbackslash{}ref as a standing input; the deduction the counted proof has to make is carried out inside the counted span, and no part of that deduction is delegated to this pointer.\par\medskip
\ifdefined\hypertarget\hypertarget{sourcepointer-6}{}\fi
\noindent\textit{Source reference.} Original reference 054K: focused-seeds/geometry/S9-morph-jumps-weighting-level-sets/ref-tags.txt, source lines 6--6. Navigation pointer only: the counted proof cites this statement by \textbackslash{}ref as a standing input; the deduction the counted proof has to make is carried out inside the counted span, and no part of that deduction is delegated to this pointer.\par\medskip
\ifdefined\hypertarget\hypertarget{sourcepointer-7}{}\fi
\noindent\textit{Source reference.} Original reference 054B: focused-seeds/geometry/S9-morph-jumps-weighting-level-sets/ref-tags.txt, source lines 7--7. Navigation pointer only: the counted proof cites this statement by \textbackslash{}ref as a standing input; the deduction the counted proof has to make is carried out inside the counted span, and no part of that deduction is delegated to this pointer.\par\medskip
\ifdefined\hypertarget\hypertarget{sourcepointer-8}{}\fi
\noindent\textit{Source reference.} Original reference 04ZC: focused-seeds/geometry/S9-morph-jumps-weighting-level-sets/ref-tags.txt, source lines 8--8. Navigation pointer only: the counted proof cites this statement by \textbackslash{}ref as a standing input; the deduction the counted proof has to make is carried out inside the counted span, and no part of that deduction is delegated to this pointer.\par\medskip
\begin{lemma}
\label{lemma-jumps-w}
Let $f : X \to Y$ be a locally quasi-finite morphism.
Let $w : X \to \mathbf{Z}$ be a weighting of $f$. Then
the level sets of the function $w$ are locally constructible in $X$.
\end{lemma}

\begin{proof}
In the proof below we will use Lemmas \ref{lemma-weighting-open} and
\ref{lemma-weighting-base-change} without further mention.
We will also use elementary properties of constructible subsets
of schemes and topological spaces, see
Topology, Section \ref{topology-section-constructible} and
Properties, Section \ref{properties-section-constructible}.
Using this the reader sees question is local on $X$ and $Y$;
details omitted. Hence we may assume $X$ and $Y$ are affine.
If we can find a surjective morphism $Y' \to Y$
of finite presentation such that the level sets of $w$
pull back to locally constructible subsets of $X' = Y' \times_Y X$,
then we conclude by Morphisms, Theorem \ref{morphisms-theorem-chevalley}.

\medskip\noindent
Assume $X$ and $Y$ affine. We may choose an immersion $X \to T$
where $T \to Y$ is finite, see
Lemma \ref{lemma-quasi-finite-separated-pass-through-finite}.
By Morphisms, Lemma \ref{morphisms-lemma-massage-finite}
after replacing $Y$ by $Y'$ surjective finite locally free over $Y$,
replacing $X$ by $Y' \times_Y X$ and $T$ by a scheme finite locally
free over $Y'$ containing $Y' \times_Y T$ as a closed subscheme,
we may assume $T$ is finite locally free over $Y$,
contains closed subschemes $T_i$ mapping isomorphically to $Y$
such that $T = \bigcup_{i = 1, \ldots, n} T_i$ (set theoretically).
Since $T_i \subset T$ is a constructible closed subset (as the image
of a finitely presented morphism $Y \to T$ of schemes), we see
that for $I \subset \{1, \ldots, n\}$ the intersection
$\bigcap_{i \in I} T_i$ is a constructible closed subset of $T$
and hence maps to a constructible closed subset of $Y$.

\medskip\noindent
For a disjoint union decomposition
$\{1, \ldots, n\} = I_1 \amalg \ldots \amalg I_r$ with nonempty parts
consider the subset $Y_{I_1, \ldots, I_r} \subset Y$
consisting of points $y \in Y$ such that $T_y = \{x_1, \ldots, x_r\}$
consists of exactly $r$ points with $x_j \in T_i \Leftrightarrow i \in I_j$.
By our remarks above this is a constructible partition of $Y$.
There exists an affine scheme $Y'$ of finite presentation
over $Y$ such that the image of $Y' \to Y$ is exactly
$Y_{I_1, \ldots, I_r}$, see Algebra, Lemma
\ref{algebra-lemma-constructible-is-image}.
Hence we may assume that $Y = Y_{I_1, \ldots, I_r}$ for some
disjoint union decomposition
$\{1, \ldots, n\} = I_1 \amalg \ldots \amalg I_r$.
In this case $T = T(1) \amalg \ldots \amalg T(r)$ with
$T(j) = \bigcap_{i \in I_j} T_i$
is a decomposition of $T$ into disjoint closed (and hence open) subsets.
Intersecting with the locally closed subscheme $X$ we obtain an analogous
decomposition $X = X(1) \amalg \ldots \amalg X(r)$ into open and closed parts.
The morphism $X(j) \to Y$ an immersion.
Since $w$ is a weighting, it follows that $w|_{X(j)}$
is locally constant\footnote{In fact, if $f : X \to Y$ is an
immersion and $w$ is a weighting of $f$, then $f$ restricts to an
open map on the locus where $w$ is nonzero.} and we conclude.
\end{proof}

\end{document}
